% !TEX root = ../../ctfp-print.tex

\lettrine[lhang=0.17]{现}{在我们已经了解了单子(monads)}，可以通过对偶性获得余单子(comonads)——只需反转范畴中的箭头并在反范畴中工作即可。

回想一下，在最基本层面上，单子涉及克莱斯利(Kleisli)箭头的组合：

\src{snippet01}
其中\code{m}是作为单子的函子。若使用字母\code{w}（倒置的\code{m}）表示余单子，则可定义余克莱斯利(co-Kleisli)箭头为如下类型的态射：

\src{snippet02}
余克莱斯利箭头的"鱼算符"(fish operator)类比定义为：

\src{snippet03}
为了使余克莱斯利箭头构成范畴，还需存在单位元——称为\code{extract}的恒等箭头：

\src{snippet04}
这是\code{return}的对偶概念。同样需要施加结合律以及左/右单位律。综合以上要素，我们可以在Haskell中定义余单子为：

\src{snippet05}
实践中我们会使用略有不同的基元操作，这将在稍后说明。

问题是：余单子在编程中有何应用？

\section{使用余单子编程}

让我们比较单子与余单子的特性。单子通过\code{return}提供将值装入容器的方法，但不提供访问内部值的接口。当然，实现单子的数据结构可能附带访问功能，但这属于额外特性。不存在从单子提取值的通用接口，正如\code{IO}单子就以永不暴露其内容为特点。

另一方面，余单子提供从中提取单一值的方法，但不提供插入值的手段。因此若将余单子视为容器，它始终预填充有内容，并允许你观察其内部。

正如克莱斯利箭头接受值并产生带上下文修饰的结果，余克莱斯利箭头则接受值及其完整上下文并产生结果。这是\newterm{上下文计算}(contextual computation)的具体体现。

\section{乘积余单子}

还记得读者单子（reader monad）吗？我们引入它来解决实现需要访问某些只读环境\code{e}的计算问题。这类计算可以表示为如下形式的纯函数：

\src{snippet06}
我们通过柯里化（currying）将它们转换为克莱斯利（Kleisli）箭头：

\src{snippet07}
但请注意，这些函数本身已经具有共克莱斯利（co-Kleisli）箭头的形式。让我们将参数整理成更便于操作的函子形式：

\src{snippet08}
我们可以通过让组合的箭头共享相同环境来定义组合算子：

\src{snippet09}
其中\code{extract}的实现简单地忽略环境：

\src{snippet10}
不出所料，乘积余单子（product comonad）可以执行与读者单子完全相同的计算。某种意义上说，余单子方式实现的环境更加自然——它遵循"上下文中的计算"这一理念。另一方面，单子（monads）具有\code{do}记法这种便捷的语法糖。

读者单子与乘积余单子之间的关联更为深刻，这源于读者函子是乘积函子的右伴随这一事实。尽管如此，一般来说，余单子涵盖的计算概念与单子有所不同。我们将在后续看到更多例子。

将\code{Product}余单子推广到任意乘积类型（包括元组和记录）是显而易见的。

\section{剖析组合}

继续对偶化的过程，我们可以进一步对偶化单子的bind和join操作。或者，我们可以重复研究鱼算子(fish operator)时采用的方法。后者似乎更具启发性。

关键在于认识到组合算子必须生成一个取\code{w a}并产生\code{c}的余克莱斯利箭头(co-Kleisli arrow)。产生\code{c}的唯一方式是对类型为\code{w b}的参数应用第二个函数：

\src{snippet11}
但我们如何生成可供\code{g}使用的\code{w b}值呢？我们手头有类型为\code{w a}的参数和函数\code{f :: w a -> b}。解决方案是定义bind的对偶操作，称为extend：

\src{snippet12}
利用\code{extend}我们可以实现组合：

\src{snippet13}
接下来能否剖析\code{extend}的内部结构？你可能会想，为什么不直接将函数\code{w a -> b}应用于参数\code{w a}？但很快你会发现无法将得到的\code{b}转换为\code{w b}。记住，余单子(comonad)没有提供提升值的机制。此时回想单子(monad)的对应构造过程，我们当时使用了\code{fmap}。这里唯一可行的方式是如果我们拥有类型为\code{w (w a)}的结构。若能将\code{w a}转化为\code{w (w a)}，这恰好就是\code{join}的对偶操作，我们称之为\code{duplicate}：

\src{snippet14}
因此，就像单子的定义一样，我们有三种等价的余单子定义：分别基于余克莱斯利箭头、\code{extend}或\code{duplicate}。以下是来自Haskell的\\ \code{Control.Comonad}库的标准定义：

\src{snippet15}
该定义提供了\code{extend}与\code{duplicate}之间的默认相互实现，因此你只需重写其中一个即可。

这些函数的直观理解基于以下观点：一般来说，余单子可以视为装满类型为\code{a}值的容器（乘积余单子是仅包含单个值的特例）。存在一个"当前"值的概念，可通过\code{extract}直接访问。余克莱斯利箭头执行聚焦于当前值的计算，但能访问周围所有值。以康威的生命游戏为例：每个单元格包含一个值（通常是\code{True}或\code{False}），对应该游戏的余单子是一个聚焦于"当前"单元格的网格。

那么\code{duplicate}的作用是什么？它将余单子容器\code{w a}转化为容器的容器\code{w (w a)}。每个新容器都聚焦于原容器\code{w a}中的不同\code{a}值。在生命游戏中，你会得到网格的嵌套结构——外层网格的每个单元格都包含一个聚焦于不同单元格的内层网格。

现在观察\code{extend}的行为。它接收一个余克莱斯利箭头和一个装满\code{a}值的余单子容器\code{w a}，将计算应用到所有\code{a}值上，用\code{b}替换它们。最终结果是装满\code{b}值的余单子容器。\code{extend}通过依次改变焦点来实现此功能，对每个\code{a}应用余克莱斯利箭头。在生命游戏场景中，余克莱斯利箭头会计算当前细胞的新状态，为此需要查看其上下文（可能是邻近细胞）。\code{extend}的默认实现展示了这一过程：首先调用\code{duplicate}生成所有可能的焦点，然后对每个焦点应用\code{f}。

\section{流余单子}

这种将容器焦点从一个元素转移到另一个元素的过程，最适合通过无限流(infinite stream)的例子来说明。这样的流与列表类似，只是它没有空构造器：

\src{snippet16}
它显然构成一个\code{Functor}：

\src{snippet17}
流的焦点是其首元素，因此\code{extract}的实现如下：

\src{snippet18}
\code{duplicate}会生成流的流，每个子流聚焦于不同的元素。

\src{snippet19}
第一个元素是原始流本身，第二个元素是原始流的尾部，第三个元素是尾部的尾部，依此类推直至无穷。

完整的实例如下：

\src{snippet20}
这是看待流的一种典型函数式方式。在命令式语言中，我们可能会从一个\code{advance}方法开始，该方法将流向前移动一个位置。而在Haskell中，\code{duplicate}能一次性生成所有移位后的流。Haskell的惰性求值机制使这种做法既可行又高效。当然，为了使\code{Stream}更具实用性，我们也可以实现类似\code{advance}的功能：

\src{snippet21}
但这个方法永远不会出现在余单子接口中。

如果你有数字信号处理经验，会立即发现流的余克莱斯利(co-Kleisli)箭头本质上就是一个数字滤波器，而\code{extend}会产生经过滤波的流。

作为简单示例，让我们实现移动平均滤波器。以下是计算流前\code{n}项和的函数：

\src{snippet22}
以下函数计算流前\code{n}个元素的平均值：

\src{snippet23}
部分应用的\code{average n}是一个余克莱斯利箭头，因此我们可以用\code{extend}在整个流上扩展该操作：

\src{snippet24}
结果就是连续移动的平均值序列。

流是单向一维余单子的典型示例。通过适当改造，可以轻松将其扩展为双向结构或多维结构。

\section{范畴论中的余单子}

在范畴论中定义一个余单子（comonad）是对偶性原理的一个直接应用。
如同单子（monad）的情形，我们从一个自函子 \code{T} 开始。
定义单子的两个自然变换 $\eta$ 和 $\mu$ 在余单子情形下被简单反转：
\begin{align*}
  \varepsilon & \Colon T \to I   \\
  \delta      & \Colon T \to T^2
\end{align*}
这些变换的分量对应着 \code{extract} 和 \code{duplicate} 操作。
余单子定律正是单子定律的镜像形式，这并不令人意外。

接着是从伴随关系导出单子的结论。对偶性会反转伴随关系：
左伴随变为右伴随，反之亦然。
由于复合 $R \circ L$ 定义了一个单子，
那么 $L \circ R$ 必然定义一个余单子。
伴随的余单位（counit）：
$$\varepsilon \Colon L \circ R \to I$$
确实就是我们在余单子定义中看到的 $\varepsilon$---
或者具体化后对应 Haskell 中的 \code{extract}。
同样我们可以使用伴随的单位（unit）：
$$\eta \Colon I \to R \circ L$$
在 $L \circ R$ 的中间插入一个 $R \circ L$
从而产生 $L \circ R \circ L \circ R$。
通过 $T$ 构造 $T^2$ 就定义了 $\delta$，
这就完成了余单子的全部定义。

我们之前也看到过单子本质上是一个幺半群。
这个命题的对偶就需要使用余幺半群（comonoid）的概念，
那么什么是余幺半群呢？
原始的"将单子视为单对象范畴"的定义并不能给出有趣的对偶结果。
当你反转所有自同构的方向时，得到的是另一个单子。
然而回想我们处理单子的方式，采用的是更一般的定义：
在单子范畴中将单子视为某个对象。
该构造基于两个态射：
\begin{align*}
  \mu  & \Colon m \otimes m \to m \\
  \eta & \Colon i \to m
\end{align*}
反转这些态射的方向就得到了单子范畴中的余幺半群：
\begin{align*}
  \delta      & \Colon m \to m \otimes m \\
  \varepsilon & \Colon m \to i
\end{align*}
可以在 Haskell 中编写余幺半群的定义：

\src{snippet25}
不过这个定义相当平凡。\code{destroy} 显然忽略其参数。

\src{snippet26}
\code{split} 只是一个函数对：

\src{snippet27}
现在考虑对偶于单子单位律的余幺半群定律。

\src{snippet28}
这里 \code{lambda} 和 \code{rho} 分别是左单位子和右单位子
（参见\hyperref[monads-categorically]{单子范畴}的定义）。
代入具体定义后得到：

\src{snippet29}
这证明了 \code{g = id}。同理第二个定律展开后得到 \code{f = id}。
综上所述：

\src{snippet30}
这表明在 Haskell 中（更一般地说，在范畴 $\Set$ 中），
每个对象都是平凡的余幺半群。

幸运的是还有其他更有趣的单子范畴可用于定义余幺半群。
其中一个就是自函子范畴。
事实证明，正如单子是自函子范畴中的幺半群，

\begin{quote}
  余单子是自函子范畴中的余幺半群。
\end{quote}

\section{存储余单子}

另一个重要的余单子例子是状态单子的对偶。
它被称为余状态余单子，或者更常见地称为存储余单子。

我们之前已见过状态单子由定义指数的伴随生成：
\begin{align*}
  L z & = z\times{}s      \\
  R a & = s \Rightarrow a
\end{align*}
我们将用相同的伴随来定义余状态余单子。余单子
由组合 $L \circ R$ 定义：
$$L (R a) = (s \Rightarrow a)\times{}s$$
将其转换为Haskell时，我们从左侧的
\code{乘积(Product)}函子与右侧的\code{读者(Reader)}函子之间的伴随开始。将\code{Reader}后接\code{Product}组合就相当于以下定义：

\src{snippet31}
在对象$a$处取伴随的上单位（counit）
是如下形态：
$$\varepsilon_a \Colon ((s \Rightarrow a)\times{}s) \to a$$
在Haskell表示法中则为：

\src{snippet32}
这将成为我们的\code{extract}方法：

\src{snippet33}
伴随的单位(unit)：

\src{snippet34}
可重写为部分应用的数据构造器：

\src{snippet35}
我们通过水平组合构建$\delta$（即\code{duplicate}）：
\begin{align*}
  \delta & \Colon L \circ R \to L \circ R \circ L \circ R \\
  \delta & = L \circ \eta \circ R
\end{align*}
我们需要将$\eta$插入到最左侧的$L$（即\code{Product}函子）中。这意味着对$\eta$（即\code{Store f}）作用于配对的左侧分量（这正是\code{Product}的\code{fmap}所做操作）。最终得到：

\src{snippet36}
（请注意，在$\delta$的公式中，$L$和$R$代表的是恒等自然变换，其分量都是恒等态射。）

以下是\code{Store}余单子的完整定义：

\src{snippet37}
你可以将\code{Store}中的\code{Reader}部分视为一个广义容器，其中元素类型\code{s}作为键索引访问\code{a}类型的值。例如，当\code{s}为\code{Int}时，\code{Reader Int a}就是无限双向流形式的\code{a}序列。\code{Store}将此容器与键类型的一个值配对。例如，\code{Reader Int a}与一个\code{Int}配对。在这种情况下，\code{extract}使用该整数索引访问无限流。你可以将\code{Store}的第二个分量视为当前的位置指针。

继续这个例子，\code{duplicate}会创建一个新的以\code{Int}为索引的无限流，这个流的元素本身就是流。特别地，在当前位置包含原始流。但当你使用其他正/负整数作为键时，会得到相对于新索引偏移后的流。

一般而言，你可以验证当\code{extract}作用于\code{duplicate}d的\code{Store}时会产生原始\code{Store}（实际上，余单子的单位律要求\code{extract . duplicate = id}）。

存储余单子作为Lens库的理论基础具有重要地位。从概念上讲，
\code{Store s a}余单子封装了使用类型\code{s}作为索引聚焦（如同镜头）数据类型\code{a}特定子结构的思想。特别地，形如：

\src{snippet38}
的函数等价于一对函数：

\src{snippet39}
如果\code{a}是乘积类型，\code{set}可实现为设置\code{a}内部类型为\code{s}的字段并返回修改后的\code{a}。类似地，\code{get}可实现为从\code{a}读取\code{s}字段的值。我们将在下一节深入探讨这些概念。

\section{挑战}

\begin{enumerate}
  \tightlist
  \item
        使用 \code{Store} 余单子实现康威生命游戏。
        提示：你为 \code{s} 选择什么类型？
\end{enumerate}